\documentclass[11pt]{article}
\usepackage[margin=2.2cm]{geometry}
\usepackage{amsmath,amssymb}
\usepackage{graphicx}
\usepackage{listings}
\usepackage{xcolor}
\usepackage{enumitem}
\usepackage{parskip}
\usepackage{booktabs}
\usepackage[hidelinks]{hyperref}

\graphicspath{{figures/}{../paper/figures/}}
\lstset{basicstyle=\ttfamily\small, columns=fullflexible, keepspaces=true, frame=single,
  rulecolor=\color{gray}, breaklines=true, showstringspaces=false, xleftmargin=4pt,
  xrightmargin=4pt, language=Python, commentstyle=\color{gray}}
\newcommand{\note}[1]{\textit{\small[note: #1]}}
\setlist{nosep, leftmargin=1.4em}

\begin{document}

{\large\textbf{Expiry-day microstructure and the settlement VWAP}}\\
Working notes: methodology, code and results, with the follow-ups from the meeting of 21 Oct\\
\small{(condensed version of the draft report, for discussion; numbers are from the current run)}
\normalsize

\hrule
\vspace{1em}

\section*{0. What changed since the last meeting}

Action items from the minutes and where they are in these notes:
\begin{enumerate}
  \item \textbf{Reasons behind each result}: a ``why'' paragraph under each hypothesis
  (\S4--\S7), with a test behind it where we could build one.
  \item \textbf{Next-morning fall in Figure 2: news after the session?} New tests in \S8. Short
  answer: mostly yes. Once we adjust for market news and drop stock-days with signs of news,
  the extra reversal on expiry days halves and is no longer significant.
  \item \textbf{BH2 plots of spread and depth}: new figures in \S5, through the whole day and
  minute by minute through the window.
  \item \textbf{Futures}: new in \S9. We rebuilt the order book of the expiring stock futures
  from the F\&O order files, and looked at the futures trades.
  \item \textbf{Detailed report}: everything here has also been written into the full report
  (new subsections in BH1--BH4, updated abstract and conclusion).
\end{enumerate}

\section*{1. The question and the four hypotheses}

Stock and index derivatives on NSE settle at the volume-weighted average price of the last
30 minutes (15:00--15:30). If someone can profit from moving that average, the window on expiry
days should look different. Following your advice we organised everything around four broad
hypotheses, each with one primary measure (expiry day vs.\ the trading day before it, same
stock, same month):

\begin{center}
\small
\begin{tabular}{llll}
\toprule
 & Hypothesis & Primary measure & Result \\
\midrule
BH1 & Settlement price is pushed & $|$settle VWAP $-$ mid at 15:00$|$ & not rejected ($p=0.21$) \\
BH2 & Window is costlier to trade in & quoted spread & rejected ($p<0.001$, $d=0.16$) \\
BH3 & One-sided pressure in window & $|$OIB$|$, OIB $=(BI-SI)/(BI+SI)$ & not rejected ($p=0.54$) \\
BH4 & Order behaviour changes & orders cancelled within 1\,s & rejected ($p<0.001$, $d=0.13$) \\
\bottomrule
\end{tabular}
\end{center}
$p$-values are paired tests across stock-months, with false discovery rate control over the four
primary tests. Each primary has supporting measures (about 30 in all) that we list in the full
report.

\section*{2. Data}

\begin{itemize}
  \item NSE order and trade files, equity cash segment, all of 2022: 242 of 243 sessions (the
  trade file for 2 Aug 2022 is corrupt).
  \item \textbf{Order book sample}: 24 sessions, the 12 monthly expiries and the trading day
  before each. For these we replay the order file and rebuild the full limit order book
  (5-second snapshots, 20 levels).
  \item \textbf{Index file}: Nifty 50 and Nifty Next 50 every second, all 243 sessions.
  \item \textbf{F\&O files}: trades and orders for March--June 2022 only (4 monthly expiries).
  \item Trade direction: each trade carries both order numbers; numbers are assigned on arrival,
  so the later one is the aggressor. No Lee--Ready needed.
\end{itemize}

\textbf{Universe} (150 stocks): 50 liquid F\&O stocks, 50 illiquid F\&O stocks, and 50
\emph{placebo} stocks with no derivatives, matched on turnover and price. The placebo group
trades on the same days but has nothing settling, so anything that also shows up there is a
property of the day, not of the settlement.

\section*{3. Pipeline}

\begin{lstlisting}[language={}]
raw NSE files  CASH_Orders / CASH_Trades / CASH_Index / FAO_Trades / FAO_Orders (.DAT.gz)
   |  nsetick (our Rust parser): fixed-width -> parquet, partitioned by date and symbol
   v
parsed/cash_trades, parsed/cash_orders, parsed/fao_trades
   |                                   |
   |  daily_tape.py: one row per       |  nsetick.build_books: replay every order,
   |  stock-session (242 sessions)     |  5-second L2 snapshots (24 sessions)
   v                                   v
tape/session=*.parquet            clob_snapshots/<symbol>/date=*/
   |  settlement_tests.py              |  s1..s8 (order book measures)
   v                                   v
s9_panel.parquet (36,300 rows)    paired expiry-vs-control tests
   |
   +--> fixed-effects regressions, randomization inference, figures, tables
\end{lstlisting}
Everything heavy runs on the HPC:
\begin{lstlisting}[language=bash]
ssh hpc
source ~/sandbox/tools/pc_env.sh            # venv + paths to the raw data
python run_all.py --stage tape --jobs 8     # 242 sessions -> tape
python run_all.py --stage report            # tests, tables, figures
python -m stage9_tape.overnight_tests       # new: section 8
python -m stage9_tape.futures_book          # new: section 9 (futures order books)
python -m stage9_tape.followups             # new: figures in sections 5 and 9
\end{lstlisting}

\textbf{Core of the tape reduction} (DuckDB over the parsed trades, one row per stock-session):
\begin{lstlisting}[language=SQL]
CASE WHEN buy_order_number > sell_order_number THEN 1.0 ELSE -1.0 END AS sgn,  -- aggressor
...
SUM(px*q) FILTER (WHERE tt BETWEEN '15:00' AND '15:30')
  / SUM(q) FILTER (WHERE tt BETWEEN '15:00' AND '15:30')        AS settle_vwap,
SUM(q) FILTER (WHERE sgn > 0 AND <5-minute block k>)            AS w{k}_bi,
SUM(q) FILTER (WHERE sgn < 0 AND <5-minute block k>)            AS w{k}_si,
\end{lstlisting}
\textbf{Estimation}: two-way fixed effects (stock and session), standard errors clustered by
session, so a market-wide move on one day can't masquerade as an expiry effect:
\begin{lstlisting}
fit = fe_ols(panel, y="next_rev",
             x=["drift", "drift_x_monthly_expiry", ..., "monthly_expiry", ...],
             fe=("symbol", "session"), cluster="session")
\end{lstlisting}
Randomization inference: re-run with the 12 expiries replaced by a random ordinary session in
the same month, 1{,}000 times.

\section*{4. BH1: is the settlement price pushed?}

\textbf{Result.} For stocks, no. The settlement VWAP is 36.0 bps from the 15:00 mid on expiry
days vs.\ 34.5 on the day before ($p=0.21$). Volume is drawn into the window on expiry days
(window share of the day's volume $+2.4$ pp more than in placebo stocks, $p<0.001$), but the
price doesn't move further.

\textbf{For the index, yes.} Nifty 50 relative to Nifty Next 50 (no derivatives) falls 6.2 bps
more over the window on index expiries (about 11 index points; randomization $p=0.008$), on
weekly and monthly expiries alike.

\begin{center}
\includegraphics[width=0.85\textwidth]{index_expiries.pdf}\\
\small Window move of Nifty 50 minus Nifty Next 50 on each index expiry of 2022.
\end{center}

\textbf{Why (reason).} The incentive to push differs by how the contract settles:
\begin{itemize}
  \item Stock F\&O have been \emph{physically settled} since 2019. A long future pays $F$ for the
  shares whatever the settlement price is, so pushing the VWAP gains nothing. Test: in the last
  minutes of the window the expiring future tracks the live cash price, not the (already almost
  fixed) VWAP. Slope of futures on cash, both relative to settle, at 15:28: \textbf{0.91}
  (se 0.06). A future settled against the VWAP would give 0.
  \item No pull of the settlement price towards option strikes either (distance to nearest
  strike unchanged, $-0.01$ intervals, se 0.01).
  \item The index is \emph{cash-settled} against the closing value, so there the incentive is
  real, and that is where we find the effect. We can't see positions, so we can't separate
  deliberate pushing from option writers unwinding cash hedges at the close.
\end{itemize}

\section*{5. BH2: is the window costlier to trade in?}

\textbf{Result.} Yes, and only in stocks with derivatives. Quoted spread over the window rises
from 5.31 to 5.60 bps ($p<0.001$); relative to placebo $+0.28$ bps ($p=0.016$). Displayed
notional in the book falls from 2.39 to 2.17 crore.

\textbf{New plots (action item 3).} Median across stocks, 5-minute smoothing, window shaded:
\begin{center}
\includegraphics[width=\textwidth]{bh2_spread_day.pdf}\\
\includegraphics[width=\textwidth]{bh2_depth_day.pdf}\\
\small Quoted spread (top) and displayed depth within 5 levels (bottom) through the session,
expiry day vs.\ day before (log scale).
\end{center}

\begin{center}
\includegraphics[width=0.95\textwidth]{bh2_window_zoom.pdf}\\
\small The last hour minute by minute. F\&O stocks: spread above and depth below the day before
throughout the window. Placebo stocks: the two days overlap.
\end{center}

Paired version: expiry minus day before, same stock and month, in log points (\%), every
15 minutes:
\begin{center}
\includegraphics[width=\textwidth]{bh2_ratio_by_time.pdf}
\end{center}

\begin{center}
\small
\begin{tabular}{llrrr}
\toprule
 & & Spread & Depth & Volatility (RV) \\
\midrule
09:30--15:00 & Liquid F\&O & $+4.5$*** & $-1.0$ & $+10.0$*** \\
             & Illiquid F\&O & $+1.7$** & $-0.6$ & $+13.8$*** \\
             & Placebo & $+1.6$* & $+1.0$ & $+11.3$*** \\
15:00--15:30 & Liquid F\&O & $+9.8$*** & $-7.2$*** & $+45.0$*** \\
             & Illiquid F\&O & $+10.3$*** & $-4.8$** & $+46.5$*** \\
             & Placebo & $+1.2$ & $-0.5$ & $+6.7$ \\
\bottomrule
\end{tabular}\\
Expiry vs.\ day before, \% (100 $\times$ log ratio), 600 stock-months per group.
\end{center}

So the spread creeps up during the expiry day in F\&O stocks and jumps in the window, where
depth also falls. Placebo stocks barely move in the window.

\textbf{Why (reason): volatility.} Realised volatility of the mid in the window is 45\% higher
on expiry days in F\&O stocks and only 7\% (n.s.) in placebo stocks. Liquidity providers widen
and thin out when prices move more (inventory and adverse-selection risk). Test: regress each
stock-month's change in log spread on its change in log volatility:
\begin{lstlisting}
d = book_paired(book, 345, 375)                 # window, expiry minus day before
X = np.column_stack([np.ones(len(d)), d.d_rv])  # intercept = change at unchanged volatility
beta = np.linalg.lstsq(X, d.d_spread, rcond=None)[0]
\end{lstlisting}
\begin{center}
\small
\begin{tabular}{lrrr}
\toprule
F\&O stocks, window & raw change & slope on $\Delta$RV & change at unchanged RV \\
\midrule
Spread & $+10.0$\%*** & 0.15 (0.01) & $+3.0$\%*** \\
Depth  & $-6.0$\%*** & $-0.10$ (0.01) & $-1.4$\% (n.s.) \\
\bottomrule
\end{tabular}
\end{center}
About 70\% of the widening and all of the fall in depth are what the higher volatility alone
would predict. The volatility itself plausibly comes from the extra (two-sided) trading drawn into
the window by expiry: institutions closing, rolling or benchmarking to the close (BH4 participation).
\note{a residual 3\% widening remains; could be market makers in the cash market also making
markets in the expiring options and pulling back there.}

\section*{6. BH3: one-sided pressure?}

\textbf{Result.} No. $|$OIB$|$ over the window is 0.12 on ordinary days in liquid stocks and
rises by only 0.02 on expiries (n.s.; paired primary test $p=0.54$). But imbalance becomes
\emph{more persistent} block to block (5-min AR coefficient 0.34 on ordinary days, $+0.15$ on
expiries, se 0.04), i.e.\ large orders worked through the window.
\begin{center}
\includegraphics[width=\textwidth]{oib.pdf}\\
\small Left: $|$OIB$|$ per 5-minute block. Right: persistence of OIB from one block to the next.
\end{center}

\textbf{Why (reason).}
\begin{itemize}
  \item Physical settlement: nobody gains from pushing (BH1), so flow at the close is the
  closing and benchmarking of positions, which comes from both sides.
  \item Positions are mostly rolled \emph{before} the window. In the futures trades (\S9) the
  next-month contract already has 53\% of stock futures volume the day before expiry and 59\%
  on expiry day, so little is left to unwind at 15:00.
  \item Participation: custodian (institutional) participation in window volume rises from
  59\% to 65\% in liquid stocks. Institutions trade towards the close price, on both sides.
\end{itemize}

\section*{7. BH4: order behaviour}

\textbf{Result.} Share of orders cancelled within one second rises from 28.6\% to 30.1\%
($p<0.001$), 2.4 pp more than in placebo stocks. Entries rise as fast as cancellations, so it is
more quoting overall rather than ``spoofing''-type place-and-pull. Concealment and IOC usage do
not change.

\textbf{Why (reason): again volatility.} Same test as in BH2, with the change in the fleeting
share on the left:
\begin{center}
\small
\begin{tabular}{lrrr}
\toprule
 & raw change (pp) & slope on $\Delta$RV & change at unchanged RV (pp) \\
\midrule
All stocks & $+1.6$ & 0.037 (0.003) & $+0.4$ (n.s.) \\
Liquid F\&O & $+2.7$ & 0.034 (0.006) & $+1.2$ (se 0.6) \\
\bottomrule
\end{tabular}
\end{center}
Algorithms refresh quotes faster when the price moves faster. Most of the rise in fleeting
orders goes with the rise in volatility.

\section*{8. Figure 2: is the next-morning fall news? (action item 2)}

Figure 2 of the draft: the window's move, oriented so that it is positive, and where the price
is the next morning (first 30 min VWAP). On monthly expiries part of the move is given back; on
ordinary days it isn't.
\begin{center}
\includegraphics[width=\textwidth]{window_path.pdf}
\end{center}
Regression version: next-morning return on the window move, two-way FE; the extra coefficient on
monthly expiries is $-0.26$ (se 0.13) in F\&O stocks and $-0.23$ (se 0.09) in placebo stocks.
Already a hint: it happens in stocks where nothing settles.

\textbf{How we tested for news.} We don't have an announcements feed, so we used what the trade
data can show:
\begin{lstlisting}
# 1. stock-days with signs of news the next morning
abnormal = next_volume / median_volume(symbol) > 2
big_gap  = |log(next_open / settle)| > 95th pct of the stock's own gaps
news = abnormal | big_gap                      # 17.5% of normal, 17.9% of expiry days

# 2. market news hits high-beta stocks harder: take it out
beta = cov(stock close-to-close, market) / var(market)     # over 2022
next_rev_beta = next_rev - beta * market_next_rev(session)

# 3. one expiry at a time (does a handful of dates carry it?)
\end{lstlisting}

\begin{center}
\includegraphics[width=0.78\textwidth]{overnight_robustness.pdf}\\
\small Extra reversal on monthly expiries under each check, 95\% intervals.
\end{center}

\begin{center}
\small
\begin{tabular}{lrr}
\toprule
Extra reversal on monthly expiries & F\&O stocks & Placebo \\
\midrule
Baseline & $-0.26$ ($p=0.04$) & $-0.23$ ($p=0.007$) \\
Drop stock-days with news signs & $-0.22$ ($p=0.08$) & $-0.09$ ($p=0.06$) \\
Adjust for market exposure (beta) & $-0.15$ ($p=0.20$) & $-0.15$ ($p=0.006$) \\
Drop 24 Feb 2022 & $-0.20$ ($p=0.14$) & $-0.17$ ($p=0.01$) \\
Beta-adjusted and no news days & $-0.13$ ($p=0.34$) & $-0.06$ ($p=0.17$) \\
\bottomrule
\end{tabular}
\end{center}

\begin{center}
\includegraphics[width=\textwidth]{overnight_by_expiry.pdf}\\
\small Slope for each monthly expiry (orange) against the 177 ordinary sessions (grey).
\end{center}

\textbf{What this says.}
\begin{itemize}
  \item It is carried by a few dates, not all expiries. 24 Feb 2022 (the day Russia invaded
  Ukraine: market crashed during the day and bounced the next morning), 31 Mar, 25 Aug and
  29 Sep stand out; 10 of the 12 expiries lie inside the middle 90\% of ordinary sessions.
  Testing the 12-expiry average against random ordinary sessions: $p=0.16$.
  \item About half of it is market news landing unevenly: high-beta stocks moved most in the
  window and moved back most when the market recovered overnight.
  \item With both adjustments what is left is small and insignificant, in both groups.
  \item It is as strong in placebo stocks, which have no settlement.
\end{itemize}
So the next-morning fall is best read as information arriving after the session (mostly
market-wide), not as the settlement price being released. \note{caveat: the earlier part of the
day (open to 15:00) is reversed much less than the window ($-0.05$), so not everything is news;
an announcements dataset (NSE corporate filings) would let us check stock news directly. Next
step.}

\section*{9. Futures (action item 4)}

\textbf{9.1 Futures order book.} F\&O order files are 99 GB a day (gz), split by stream; stock
futures sit in 3 of the 12--15 parts. We parse only those, only the expiring contract, no spread
orders, and replay each future's book exactly as for cash:
\begin{lstlisting}
nsetick.parse(input=part, layout="fao_orders", partition_by="symbol",
    where="instrument == 'FUTSTK' and spread_type == '*' "
          "and expiry_date == '28Apr2022' and symbol in (...)")
nsetick.build_books(input=orders_dir, out=books_dir, interval_secs=5.0, levels=5)
\end{lstlisting}
4 expiries (Mar--Jun) and the day before each, same 100 F\&O stocks.

\note{running on the HPC now: the order files are about 36 GB a day for the stock-futures
parts alone and the NSE share reads at about 30 MB/s, so the 8 sessions take a few hours.
Results go here.}

\textbf{9.2 Futures trades.} Expiring (near) and next-month contract, expiry day vs.\ day
before, 100 stocks $\times$ 4 months:
\begin{center}
\small
\begin{tabular}{llrrr}
\toprule
 & & Day before & Expiry & Difference \\
\midrule
Near (expiring) & share of the stock's futures volume (\%) & 47.4 & 41.2 & $-6.3$*** \\
                & window share of the day's volume (\%) & 12.8 & 11.1 & $-1.7$*** \\
                & $|$OIB$|$ in window & 0.18 & 0.18 & $-0.01$ \\
Next month      & share of futures volume (\%) & 52.8 & 59.1 & $+6.3$*** \\
                & $|$OIB$|$ in window & 0.15 & 0.13 & $-0.02$*** \\
\bottomrule
\end{tabular}
\end{center}
The roll is mostly done before the expiry day; trading in the expiring future does
\emph{not} concentrate in the window, and there is no one-sided pressure in it. This is
consistent with BH1/BH3 for stocks.

\section*{10. Summary and next steps}
\begin{itemize}
  \item BH1: stock settlement not pushed (physical settlement removes the incentive); index
  settlement is ($-6.2$ bps vs.\ Next 50).
  \item BH2: window costlier in F\&O stocks only; mostly explained by higher volatility in
  the window.
  \item BH3: no one-sided pressure; flow is two-sided closing, positions rolled earlier.
  \item BH4: more fleeting orders, mostly explained by volatility.
  \item Figure 2: next-morning fall mostly overnight information, not settlement.
\end{itemize}
Next: (i) corporate announcements data to test stock news directly; (ii) futures book for more
sessions (not only expiry and day before); (iii) finish the futures order book comparison (\S9.1).

\end{document}
